\section{Methodology}
\label{sec:methodology}

We propose \textbf{OpenDC-Infer-LC}, an open, hardware-neutral benchmark for
long-context and retrieval-augmented generation (RAG) inference serving in
data-center environments. It complements model-centric long-context evaluations
and generic throughput benchmarks: rather than reporting peak tokens per second,
OpenDC-Infer-LC measures whether a serving \emph{system} can sustain long-context
workloads under realistic service-level objectives (SLOs), and additionally
reports cache behavior, multi-node capacity, reliability, and energy efficiency.
The main body covers the objective and system under test
(Section~\ref{subsec:objective}), the benchmark definition---workloads, models,
and SLOs (Section~\ref{subsec:definition})---and the scoring metrics
(Section~\ref{subsec:metrics}); the full submission protocol (endpoint contract,
cache-state taxonomy, harness procedure, leaderboard rules, and corpus mechanics)
is deferred to Appendix~\ref{sec:protocol_appendix}.

\subsection{Objective and System Boundary}
\label{subsec:objective}
The central question is:
\vspace{-0.25\baselineskip}
\begin{quote}\itshape
``Given a fixed model, tokenizer, workload, and latency requirement, how many
long-context or RAG requests can a cluster serving system process reliably?''
\end{quote}
\vspace{-0.25\baselineskip}

This differs from endpoint-only benchmarking. OpenDC-Infer-LC evaluates the full
serving system---scheduler, container runtime, accelerator allocation,
interconnect, storage and cache paths, telemetry, serving backend, and request
load generator---to measure production-relevant capacity rather than whether an
endpoint is merely reachable. We target long-context and RAG-style inference
because these workloads dominate a growing share of production traffic (enterprise
document QA, codebase reasoning, agent memory, multi-turn retrieval-augmented
assistants) and, relative to short chat, place far greater pressure on prefill
efficiency, KV-cache memory, request scheduling, prefix reuse, and multi-node
scaling. The benchmark is hardware- and backend-neutral: any accelerator (AMD,
NVIDIA, TPU, custom) and any serving backend (vLLM, SGLang, TensorRT-LLM, TGI, or a
proprietary runtime) may participate, provided it follows the measurement and
submission rules.

\subsection{Benchmark Definition}
\label{subsec:definition}
A benchmark track fixes three things: \emph{what} is served (the workloads and
corpus), with \emph{which} model and tokenizer, and to \emph{what} latency target
(the SLO profile). All three are frozen so that submissions differ only in the
system under test.

\subsubsection{Workloads and Corpus}
\label{subsec:workloads}
OpenDC-Infer-LC defines six workloads (Table~\ref{tab:workloads}) chosen to span
the shapes that dominate production long-context and RAG traffic. Four
\emph{length/cache} workloads isolate context scaling and prefix reuse: for LC-8K,
LC-32K, and LC-128K a request is a single long context, a user query, and a fixed
maximum generation length; LC-Cache instead shares one long prefix across many
requests with differing queries, exercising prefix reuse and KV-cache management.
Two \emph{retrieval-shape} workloads capture how enterprise RAG actually arrives.
\textbf{RAG-TopK} models vector-database retrieval: each request concatenates $k{=}64$
independently retrieved passages ($\sim$512 tokens each) from \emph{different}
documents, exactly one of which carries the answer while the rest are hard
distractors, and the retrieved set varies per query---so unlike the single-haystack
workloads it stresses \emph{fragmented} prefill and \emph{low cross-query cache
reuse}. \textbf{RAG-Report} models report/summary generation: a multi-document
context with a \emph{long structured output} (2048 tokens) that must reproduce a set
of required facts, exercising the \emph{decode-bound} regime and long-generation
tail latency that the short-output workloads never reach.
Because these synthetic workloads trade external validity for control, we add two
\emph{public real-text} retrieval workloads for external validity:
\textbf{RAG-RealQA} (real Wikipedia multi-document QA from LongBench HotpotQA) and
\textbf{RAG-SciQA} (scientific-paper QA from LongBench QASPER). Each real example is
packed to the same $32$K budget with additional real passages drawn from the corpus
and graded by token-F1 (the standard multi-document-QA metric); they are the
recognized-benchmark analogue of RAG-TopK and confirm the synthetic findings carry to
real text and a non-Wikipedia domain (\S\ref{subsec:rag_details}).

\begin{table*}[t]
\centering
\small
\caption{OpenDC-Infer-LC workload definitions.}
\label{tab:workloads}
\begin{tabular}{lccc}
\toprule
\textbf{Workload} & \textbf{Input Tokens} & \textbf{Max Output Tokens} & \textbf{Target Scenario} \\
\midrule
LC-8K & 8K & 512 & Standard RAG \\
LC-32K & 32K & 512 & Enterprise document QA \\
LC-128K & 128K & 1024 & Long-document and codebase QA \\
LC-Cache & 32K shared prefix $+$ multi-turn queries & 256 per turn & Agentic and repeated-prefix RAG \\
\midrule
RAG-TopK & $k{=}64$ retrieved passages ($\sim$512 tok each, $\approx$32K) & 256 & Vector-DB top-$k$ retrieval \\
RAG-Report & 32K multi-document & 2048 & Report / summary generation \\
\bottomrule
\end{tabular}
\end{table*}

The corpus is \textbf{fully synthetic and procedurally generated from a seeded
random number generator}---never sourced from public documents---using a
RULER-style needle-in-a-haystack template with ``magic number'' needles as
exact-match answer keys (for RAG-TopK the needle sits in one retrieved passage and
the others are distractors; for RAG-Report the long output must reproduce a set of
required needles). Synthetic generation gives metric integrity (a repeated real
passage would inflate prefix-cache hit rates and distort the cache metrics), license
cleanliness, and bit-for-bit reproducibility; each prompt is fitted to its target
token length \emph{exactly} ($\leq8$-token drift) per tokenizer and pinned by a
\texttt{dataset\_version\_hash}. Full generation details are in
Appendix~\ref{subsec:corpus_details}.

\subsubsection{Models and Tokenizers}
\label{subsec:model_selection}
Each track fixes the model checkpoint, tokenizer, maximum output length, and
sampling parameters. Submissions may vary hardware, serving framework, parallelism,
quantization, and cache optimizations, but must use the same benchmark-visible
model and tokenizer. We define two open-weight mixture-of-experts (MoE) tiers:
\begin{itemize}
    \item \textbf{Medium tier: Qwen3.5-35B-A3B} (35\,B total / 3\,B active),
    intended for single-node and small-cluster submissions.
    \item \textbf{Large tier: Qwen3-235B-A22B-Instruct-2507} (235\,B total /
    22\,B active, 256K native context), the main data-center leaderboard model.
\end{itemize}
Sampling is deterministic ($\texttt{temperature}{=}0$, $\texttt{top\_p}{=}1$). The
runner re-tokenizes prompts and outputs with the declared tokenizer unless a
submission provides a validated server-side token-accounting contract.

\subsubsection{Service-Level Objectives}
\label{subsec:slo}
The benchmark defines three SLO profiles (Table~\ref{tab:slo}). The
\emph{interactive} profile applies to LC-8K, LC-32K, and RAG-TopK; the \emph{heavy}
profile applies to LC-128K; and a \emph{generation} profile---which relaxes the
end-to-end bound to admit long outputs while keeping the per-token (TPOT)
bound---applies to the long-output RAG-Report. LC-Cache is reported under the
profile matching its configured context length and cache state. We gate on p95
rather than mean because production inference is tail-bound.

\begin{table*}[t]
\centering
\caption{Service-level objectives used in OpenDC-Infer-LC v1.0.}
\label{tab:slo}
\begin{tabular}{lcccc}
\toprule
\textbf{SLO Profile} & \textbf{TTFT p95} & \textbf{TPOT p95} & \textbf{E2E p95} & \textbf{Success Rate} \\
\midrule
Interactive RAG & $<3$\,s & $<120$\,ms/token & $<60$\,s & $>99\%$ \\
Heavy Long-Context & $<10$\,s & $<200$\,ms/token & $<180$\,s & $>98\%$ \\
Generation (long-output) & $<10$\,s & $<200$\,ms/token & $<300$\,s & $>98\%$ \\
\bottomrule
\end{tabular}
\end{table*}

\subsection{Metrics}
\label{subsec:metrics}
A run is scored on one primary metric (SLO-qualified goodput), a broad set of
secondary diagnostics (Table~\ref{tab:metrics}; full list in
Appendix~\ref{subsec:secondary_full}), and a correctness guardrail.

\subsubsection{Primary Metric: SLO-Qualified Goodput}
\label{subsec:goodput}
The primary metric is \textbf{SLO-qualified output-token goodput}. A run is
\emph{SLO-qualified} if its aggregate p95 TTFT, p95 TPOT, p95 E2E, and success rate
all satisfy the selected profile over the full measurement window. Let
$\mathcal{R}$ be the requests in a run and $y_i$ the output tokens of request $i$.
Define $v_i{=}1$ if request $i$ is successful (not timed out, cancelled, or
malformed), produces a valid final output, \emph{and} passes the quality guardrail
(Section~\ref{subsec:quality}); otherwise $v_i{=}0$. For duration $T$,
\begin{equation}
\mathrm{Goodput}_{\mathrm{SLO}}
= \frac{\sum_{i \in \mathcal{R}} v_i\, y_i}{T}.
\label{eq:goodput}
\end{equation}
The quality guardrail is a \emph{separate} gate from the latency/success
qualification: it does not enter SLO-qualification but enters the metric through
$v_i$, zeroing the token contribution of quality-failing requests. A configuration
can therefore be SLO-qualified on latency yet see its goodput collapse because it
answers incorrectly, so latency-qualification and answer correctness must be read
together. A non-qualified run is ineligible for ranking regardless of raw
throughput. For capacity search the harness raises the offered load until the SLO
is violated, reporting the highest qualified point and retaining the first
violating point as the saturation boundary. For multi-node submissions we report
scaling efficiency
\begin{equation}
\mathrm{Efficiency}(N) = \frac{G(N)}{N \cdot G(1)},
\label{eq:efficiency}
\end{equation}
where $N$ is the node count and $G$ is, by default, $\mathrm{Goodput}_{\mathrm{SLO}}$;
when a configuration fails the quality guardrail \emph{uniformly} across the sweep
(so goodput is a near-constant quality factor below the raw rate and would mask the
system-level scaling), we set $G$ to the raw SLO-qualified output throughput and
mark the result.

\subsubsection{Quality Guardrail}
\label{subsec:quality}
A lightweight quality guardrail prevents trading correctness for speed. Each prompt
carries a ground-truth answer (the needle value); the per-request quality check is
an \textbf{exact match} between the model's parsed answer and that value (for
multi-key prompts, the value of the \emph{queried} key), and LC-Cache additionally
checks answer consistency across multi-turn queries sharing the same prefix. The
per-run \emph{quality score} is the fraction of successful requests that pass, and a
submission is valid only if it stays within tolerance of a BF16 reference run:
\begin{equation}
\mathrm{QualityRatio}
= \frac{\mathrm{QualityScore}_{\mathrm{submission}}}{\mathrm{QualityScore}_{\mathrm{reference}}}
\;\geq\; 0.98 \;\;(\text{default}),
\label{eq:qualityratio}
\end{equation}
so a quantized or optimized system cannot climb the leaderboard by emitting faster
but incorrect answers.

\subsection{Load Generation}
\label{subsec:load_generation}
OpenDC-Infer-LC uses two load-generation modes. The \textbf{closed-loop} mode holds
a fixed number of concurrent users and sweeps exponentially increasing concurrency
($1,2,4,8,16,32,64,\ldots$) until the SLO is violated, yielding the capacity curve.
The \textbf{open-loop} mode issues a Poisson arrival process at a fixed offered
rate, swept upward to saturation, better reflecting production traffic where
requests arrive independently. For each workload we report the highest
SLO-qualified offered load and retain the first non-qualified point as the
(unranked) saturation boundary. The endpoint measurement contract, cache-state
rules, reference-harness procedure, and leaderboard/submission rules are specified
in Appendix~\ref{sec:protocol_appendix}.
